\section{Cross-Domain Synthesis}
\label{sec:synthesis}

The domain analyses connect agent decisions to three scientific relationships: how an action reduces uncertainty relevant to a research question, how disciplinary conditions constrain that action, and how its outcome becomes a result that can be checked or used in subsequent research. The following synthesis develops these relationships through selected primary studies and relevant methodology for inference under adaptive allocation. The cases support conditional conclusions about the evaluated tasks, rather than estimates of their prevalence across all 662 works.

\subsection{Allocation of Scientific Decisions Between Agents and Researchers}
\label{sec:workflow-delegation}

The central commonality is delegation of consequential choices within an investigation. An agent can improve candidate selection, decide which represented hypothesis survives, or change the model that makes hypotheses expressible. These choices address different uncertainties and require different authority. A useful comparison therefore asks both what the feedback changes and which scientific decisions were already made in constructing the task.

The Robot Scientist makes this distinction particularly explicit. Researchers supplied a Prolog model of yeast aromatic amino-acid metabolism and a set of possible gene--enzyme assignments. The system generated competing assignments, selected growth assays that discriminated among them, and used results to reject hypotheses and choose subsequent experiments~\cite{king2004}. Thus, both experiment selection and revision of the surviving hypothesis set were delegated. Construction of the pathway representation and specification of the biological question remained prior scientific inputs. The 2004 demonstration reconstructed known gene functions, so its significance lies in the integrated hypothesis-testing process and its efficiency, rather than the discovery of a previously unknown function.

Materials optimization delegates a different decision. In the mobile robotic chemist study, five researcher-proposed hypotheses motivated a ten-variable formulation space. Bayesian optimization selected formulations from measured hydrogen evolution; unfavorable results reduced the use of some proposed ingredients~\cite{burger2020}. The system revised experimental allocation within supplied chemical possibilities. Such a response is sufficient to improve empirical selection without establishing why an ingredient changes photocatalytic performance. The contrast with Robot Scientist concerns the purpose of information: ranking candidate performance and discriminating gene--enzyme assignments call for different experiments, even when both systems implement a design--execute--update cycle.

Control of the next action also extends beyond candidate search. In k-agents, a supplied quantum-experiment procedure is decomposed into stages; inspection of results governs parameter adjustment, retries, and progression to subsequent stages~\cite{asd0090}. In Sciencecraft, onboard analysis can introduce observation goals and planning can revise an executable mission schedule, within scientific priorities and spacecraft constraints supplied by the mission team~\cite{asd0645}. These cases delegate decisions that affect what will actually be measured. Their authority is therefore different from producing a recommendation that another actor must translate into action.

Conversely, broad participation in scientific discussion, design, coding, and writing can coexist with human control of execution and acceptance. Virtual Lab agents proposed tools and wrote scripts, while researchers executed all scripts and performed biological validation~\cite{asd0054}. SciSciGPT translates user research requests into literature, database, analysis, and visualization tasks~\cite{asd0624}. Such systems redistribute intellectual and technical work, but the evaluated workflow must determine whether a proposal is adopted and which actor controls the corresponding intervention. Producing outputs at several research stages does not by itself establish authority over those stages.

The cases also distinguish selection within a representation from changing the representation itself. In AI Theorist, spectral discrepancies prompted revisions to electronic-state assumptions and subsequent calculations~\cite{asd1081}. In the curious formulation robot study, researchers introduced temperature as an explanatory variable after inconsistent exploration results~\cite{asd0575}. The latter changed which conditions had to be recorded and controlled; it could not be achieved simply by ranking more mixtures under the original description. Reformulation can initially expand uncertainty by revealing possibilities excluded from the original model. Its value is to make those possibilities investigable. Model revision and problem reformulation thus have scientific consequences beyond more candidate proposals, and their allocation differs between systems. Workflow breadth is a poor substitute for identifying who can change the assumptions on which subsequent decisions depend.

\subsection{How Disciplinary Research Conditions Shape Delegation}
\label{sec:conditions-delegation}

The conditions in Section~\ref{sec:coverage-imbalance} change the appropriate strategy, including what to select next and when a result can be accepted. They are not merely obstacles added to an otherwise identical workflow. Their consequences follow from whether evaluation checks a formal object, samples a physical response, acquires an unrepeatable event, or measures a proxy for a more distant outcome.

For an explicitly verifiable mathematical construction, the policy can explore freely while the final result is checked against fixed conditions. FunSearch makes this separation operational through executable evaluation~\cite{asd0512}. AlphaEvolve's finite witness, considered in Section~\ref{sec:domains}, supplies a related example of search followed by a mathematical implication. Empirical engineering and chemical search instead samples variable physical outcomes~\cite{asd0504,asd0734}. Here, the best observation can reflect both design quality and measurement variation. Repeated evaluation and characterization help determine whether an apparent improvement is a reproducible property of the selected object. These are different reasons to allocate a further trial: verify a candidate's logical conditions, estimate its response, or investigate its identity.

Repeatable interventions also permit different scientific objectives. A materials optimizer can choose the formulation expected to improve performance, while a biological experiment selector can choose an assay whose alternative outcomes discriminate hypotheses~\cite{burger2020,king2004}. When instrument state is uncertain, calibration may be more informative than testing another candidate~\cite{asd0090}. Consequently, an appropriate next action depends on the uncertainty limiting the stated claim. Shorter feedback cycles help only if the resulting measurements address that uncertainty; throughput alone cannot choose between exploitation, discrimination, and recalibration.

For transient Earth or astronomical events, a missed observation may be irrecoverable. Sensorweb acquisition, telescope scheduling, and mission replanning therefore address the timing and allocation of sensing opportunities~\cite{asd0722,asd0136,asd0645}. Follow-up can add a new observable or time point even when the underlying process cannot be experimentally controlled. The relevant scientific comparison concerns what evidence the schedule makes available for an environmental or physical inference. Number of executed commands is an operational endpoint; coverage of a transient's informative phases is a different acquisition objective. The reports establish responsive acquisition, while the benefit to a particular explanation depends on the observations obtained.

Proxy-based investigations require a further strategic distinction. BioDiscoveryAgent's retrospective gene selections, Virtual Lab's measured binding phenotypes, and RAVA's image-level diagnoses concern accessible biological or visual endpoints~\cite{asd0020,asd0054,asd1133}. Such endpoints can answer well-defined intermediate questions. If the intended claim concerns therapeutic benefit or crop management, the next informative investigation must address the connection to that outcome. More precise prediction of the same proxy does not by itself test this connection. Scholarly-data analysis and institutional simulation differ for a related reason: an empirical network describes recorded interactions, whereas changing a simulated reviewer tests consequences within a specified behavioral model~\cite{asd0624,asd0786}. Evidence about human institutional responses requires a corresponding empirical comparison.

Transfer tests identify which of these relations must be re-established. NMR deployment required local recharacterization; changes in image data or training altered the relative advantage of active vision~\cite{asd1064,asd1133}. Microscopy configurations could remain executable even as their performance became harder to predict on new tasks~\cite{asd1009}. These results distinguish three targets of reuse: the procedure, its comparative advantage, and the scientific relation inferred from its outputs. Preserving the first is insufficient evidence for the other two. The appropriate transfer test follows from the claim: local calibration for an instrument strategy, outcome comparison for a selection method, or new evidence for a biological or physical relation.

\subsection{Demonstrated Scientific Gains and Their Attribution}
\label{sec:comparative-gains}

The primary comparisons establish several kinds of gain: more informative experiment selection, improved candidate performance, better search procedures, and more effective access to scientific tools. Their attribution depends on the comparator. A comparison with an initial formulation establishes improvement during a campaign; a comparison with another selection policy addresses how experiments are chosen; a comparison with the same underlying language model addresses the value of the added system. Table~\ref{tab:comparative-gains} separates these interpretations.

\begin{table}[H]
\caption{Comparative evidence for scientific gains in selected primary studies. Metrics retain their study-specific meanings and are not pooled across domains.}
\label{tab:comparative-gains}
\scriptsize
\begin{tabularx}{\textwidth}{@{}p{0.14\textwidth}>{\raggedright\arraybackslash}X>{\raggedright\arraybackslash}X>{\raggedright\arraybackslash}X@{}}
\toprule
Study & Comparator and endpoint & Reported result & Scope of attribution \\
\midrule
BEAR~\cite{asd0504} & Adaptive physical design search versus a benchmark derived from 1,800 grid samples; measured toughness & A batch campaign matched benchmark toughness after 32 experiments (12 h) and exceeded it after 64 (24 h) & Campaign-level selection efficiency within one structural family; benchmark construction is an additional cost \\
Mobile robotic chemist~\cite{burger2020} & Final versus initial formulation; hydrogen-evolution rate & 688 experiments in eight days; approximately sixfold improvement & Joint outcome of automated experimentation and adaptive search; no sixfold advantage over a matched human or random-search control \\
FunSearch~\cite{asd0512} & Evolved versus first-fit and best-fit heuristics; excess bins relative to a lower bound & On OR4: 2.47\% excess bins versus 5.23\% and 4.94\%, respectively & Better tested packing heuristics; the percentages do not measure search cost or the number of bins saved directly \\
Bio\-Discovery\-Agent~\cite{asd0020} & Claude 3.5 Sonnet without tools versus the best tabulated baseline per dataset; five-round hit ratio & CAR-T: 0.130 versus 0.069; Carnevale: 0.042 versus 0.047; improvement on five of six datasets & Retrospective retrieval of perturbation outcomes; comparison includes pretrained biological knowledge \\
ChemCrow~\cite{bran2024} & Tool-assisted system versus GPT-4 without tools; expert assessment of chemistry tasks & Expert assessments favored tool-assisted responses & Supports the system-level benefit of tool access and orchestration; does not isolate each component's contribution \\
\bottomrule
\end{tabularx}
\end{table}

Experiment-selection efficiency has evidence predating LLM-based systems. The Robot Scientist compared its cost-aware selection strategy with cheaper-first and random alternatives, evaluating the accuracy of remaining hypotheses within the supplied biological model~\cite{king2004}. BEAR provides a physical engineering comparison: adaptive campaigns identified high-toughness structures with far fewer trials than were used to construct the grid benchmark~\cite{asd0504}. The study expended more than 2,500 experiments in establishing and examining this comparison. The 32-trial result therefore characterizes the demonstrated search campaign, rather than the total cost of creating the platform and its benchmark. Together, these studies show how selecting informative experiments can improve research efficiency independently of language-model reasoning.

Improved measured performance supports a different conclusion from superiority of the selection method. The mobile robotic chemist's approximately sixfold increase in hydrogen evolution establishes successful optimization within its formulation space~\cite{burger2020}. Without a matched alternative campaign, that ratio cannot apportion the improvement among throughput, Bayesian selection, and the researcher-designed ingredient space. Similarly, Granda and colleagues combined automated reactivity screening with manual follow-up, isolation, and characterization of selected transformations~\cite{asd0734}. The scientific result depends on both the screening process and the subsequent work that establishes the chemical identity and significance of its outputs. These examples reveal a contribution distributed across investigation stages, rather than a single algorithmic gain measured by the final artifact.

FunSearch demonstrates an additional route to contribution: generating a reusable procedure that yields improved solutions~\cite{asd0512}. Its bin-packing comparison tests the resulting heuristic against established methods on separate benchmark instances. Its mathematical results also include new constructions, with further human analysis of discovered structure enabling additional improvements. The LLM supplies program variations; evolutionary selection and executable evaluation determine which variations survive. This combination expands the search over candidate procedures, while the reported heuristic advantage concerns the completed system's output. It does not imply that an unaided LLM solves the mathematical problem or that greater workflow breadth caused the improvement.

BioDiscoveryAgent addresses whether pretrained biological knowledge can improve sequential selection relative to data-driven acquisition baselines~\cite{asd0020}. Its five-round evaluation selects batches of 128 genes, reduced to 32 for the smaller dataset, and retrieves their measured outcomes from past screens. Claude 3.5 Sonnet without tools exceeds the strongest tabulated baseline on five of six datasets, with substantial variation between biological contexts (Table~\ref{tab:comparative-gains}). The early-round advantage is consistent with informative prior knowledge mitigating cold start. Because the LLM and the acquisition baselines have different prior information, the result does not isolate feedback-based reasoning from pretrained knowledge. Component comparisons further show that added tools improve the weaker Claude 3 Haiku model more consistently than Claude 3.5 Sonnet. Incorporating prior knowledge and adding orchestration are therefore distinguishable contributions whose benefits depend on the model and task.

ChemCrow's comparison with tool-free GPT-4 more directly tests an extension to the same base model~\cite{bran2024}. Expert preference supports the added system's utility for the tested chemistry tasks, while leaving retrieval, specialized algorithms, and orchestration jointly responsible. Across the cases, demonstrable gains attach to particular decisions, outputs, and resource conditions. A stronger claim about a decision paradigm requires a comparator that preserves the available information, tools, experimental budget, and human contribution relevant to that claim. The current evidence supports complementary contributions from symbolic hypothesis testing, optimization, planning, learned policies, and LLM-based integration; it does not establish a universal advantage for one architecture.

\subsection{Adaptive Investigation and the Validity of Scientific Evidence}
\label{sec:adaptive-evidence}

An agent that chooses experiments changes the evidence available for subsequent inference. Its scientific contribution therefore depends on the relation between the decision objective, the distribution of acquired observations, and the claim made from those observations. The distinction between property optimization and hypothesis discrimination in Section~\ref{sec:workflow-delegation} has an important consequence: a policy can become better at finding a desirable object while learning less about other parts of the system. Exploration within Bayesian optimization can support the search for an optimum without being designed to identify a mechanism or characterize the entire response surface.

This divergence is observable in a computational materials study. Yoshida and colleagues evaluated linked autonomous searches for magnetization, Curie temperature, and spin polarization using a fixed dataset of 16,908 alloys calculated by density functional theory. Transferring learned representations improved some search outcomes, yet the Curie-temperature search preferentially acquired high-temperature examples, and predictive performance on the remaining materials, which included a substantial fraction with low or zero Curie temperatures, temporarily deteriorated~\cite{asd0513}. Better candidate acquisition and better prediction throughout the search space were therefore separate achievements. This result concerns a computed candidate library, rather than networked physical laboratories. It also shows why the scientific value of a policy must be assessed against its intended question: locating a high-performing material and estimating a general composition--property relation require different coverage.

Mechanism identification imposes an additional constraint. If competing hypotheses predict the same observation distribution for every accessible action and history, a policy restricted to those actions cannot distinguish them. More iterations can improve precision without resolving this observational equivalence. The available interventions, measured quantities, and hypothesis representation jointly determine what is identifiable. In the Robot Scientist, growth-rescue assays were informative because different gene--enzyme assignments implied different outcomes within the supplied pathway model~\cite{king2004}. The inference remained conditional on that representation. In the AI Theorist study, proposed polarization-dependent spectra identify an additional measurement relevant to testing the excitonic assignments; the corresponding experiment had not been performed~\cite{asd1081}. Improving a fit to existing observations and acquiring an observation that discriminates between explanations thus address different scientific needs. For an agent, recognizing when another measurement or model class is required is a distinct decision from continuing the current search.

\begin{samepage}
Adaptive acquisition does not, by itself, invalidate inference. Let $D_t=\{(a_s,y_s)\}_{s=1}^{t}$ record actions and outcomes, $\pi$ denote the action-selection policy, and $\theta$ index the observation model. For a fixed number $T$ of actions, when the policy depends on $\theta$ only through the recorded history, the joint distribution factorizes as
\begin{equation}
p_\theta(D_T)=\prod_{t=1}^{T}\pi(a_t\mid D_{t-1})\,p_\theta(y_t\mid a_t,D_{t-1}).
\label{eq:adaptive-evidence}
\end{equation}
\end{samepage}
The policy factors contain no additional dependence on $\theta$, so a correctly specified likelihood built from the sequential observation factors can accommodate the adaptive design. This argument assumes a fully recorded decision process and a fixed horizon; it does not establish validity for arbitrary stopping rules, unrecorded selection, or model misspecification. Nor does it guarantee that conventional confidence intervals retain their nominal coverage. Hadad and colleagues show that adaptive allocation can produce biased or heavy-tailed conventional estimators and develop inference under explicit conditions on assignment probabilities and weighting~\cite{hadad2021}. Their results concern randomized adaptive experiments; they are not an automatic correction for deterministic laboratory optimizers. Where some actions are never selected, reweighting alone cannot reconstruct evidence about them without further assumptions or new observations.

Interpreting an acquired dataset consequently requires attention to how it was generated. In the sputter-epitaxy study, a random forest fitted to the initial optimization data achieved an in-sample $R^2$ of 0.92 but approximately 0.43 under five-fold cross-validation. The authors explicitly attributed the more limited predictive performance to sample size and the acquisition-biased distribution, and used the model as a local interpretable surrogate~\cite{asd0999}. A random split of such a dataset still tests prediction within the acquired distribution; it does not establish accuracy throughout the original design space. Feature importance and partial dependence summarize the fitted response under those conditions, without independently identifying a microscopic causal mechanism. The subsequent experiments using the inferred growth rule, discussed in Section~\ref{sec:domains}, supply a different kind of support: the rule guided new decisions beyond fitting the original acquisition data.

The preceding inference also depends on an adequate observation model. A residual between prediction and measurement does not, by itself, identify whether the predictive model, measurement procedure, or physical explanation requires revision. Section~\ref{sec:model-observation} examines how selected experimental systems address these different possibilities.

The endpoint determines how adaptive selection affects the final claim. For AlphaEvolve's mathematical construction, correctness rests on explicit geometric conditions and a valid implication from the finite witness~\cite{asd0063}. Adaptive search over many candidates does not weaken that implication once its premises are checked. An empirical optimum or association instead depends on measurement variability, the sampled conditions, and the relevant inferential design. Likewise, a low predicted improvement can justify ending a resource-limited optimization campaign without establishing a unique mechanism or excluding hypotheses outside the represented family. Scientific autonomy therefore includes deciding whether the available evidence answers the research question, whether a different experiment is needed, and whether the current formulation can answer it at all. This analysis extends workflow participation to the control of what can be learned from the investigation.

\subsection{Scientific Responses to Model--Observation Discrepancies}
\label{sec:model-observation}

Experimental feedback can change both the assessment of a candidate and the system's diagnosis of the source of uncertainty. A disagreement can concern a numerical prediction within an accepted model, the extraction of a physical quantity from a measured signal, or the assumptions defining the explanation itself. These require different subsequent investigations. BEAR's simulation-informed compression experiments, NMR strategy transfer, and the AI Theorist's optical-spectrum analysis are selected here because their reports document these distinct responses. They concern engineered structures and physical systems; their comparison does not presume that the same diagnostic operations apply to formal proofs or institutional simulations.

BEAR illustrates correction of a useful but imperfect physical prior. In crossed-barrel compression, finite-element analysis supplied predicted resilience $\widetilde U_E(x)$, and physical experiments supplied the energy $U_E(x)$ absorbed up to yield under the study's measurement convention. Rather than replace the simulation, a Gaussian process learned its discrepancy. For the different objective of toughness $U(x)$, defined as the area under the full measured force--displacement curve, the researchers instead introduced an effective-length response using predicted yield force $\widetilde F_y(x)$~\cite{asd0615}:
\begin{equation}
\delta(x)=U_E(x)-\widetilde U_E(x),\qquad
L(x)=\frac{U(x)}{\widetilde F_y(x)}.
\label{eq:bear-response}
\end{equation}
These are two study-specific response representations, not interchangeable physical objectives. Resilience and toughness were component-level energies measured in joules. The effective length expresses the displacement that would absorb energy $U$ at constant force $\widetilde F_y$; it is not a direct measurement of ductility. The elastic-response prediction did not correlate with toughness, and direct use of that prior did not accelerate simulated toughness search. Nevertheless, predicted yield force remained useful for transforming the response and downweighting designs likely to exceed the 4.8-kN test threshold~\cite{asd0615}. Thus, an approximate model can inform the scale or feasible region of a decision even when its original output poorly predicts the target. The measurement protocol helps determine which model information is useful. Researchers supplied these representations and constraints; the optimizer updated beliefs and selected designs within them. Successful discrepancy learning supports local prediction and selection without establishing a physical explanation of the residual.

NMR transfer exposes a different problem: the procedure for estimating the local physical model may cease to be appropriate. On processor Q4, two four-line multiplets were resolved, but inner transitions in the third overlapped. The original fitter required three resolved quartets. The agent adapted line assignment to infer couplings from the two resolved quartets and constrain them with the outer lines of the third; separate X- and Y-excitation spectra tested consistency~\cite{asd1064}. This preserved the three-spin model class while changing how its parameters were extracted. Importantly, a plausible fit was insufficient: none of nine initial spectral pairs passed all characterization checks, so the agent stopped and requested operator adjustment. A later acquisition passed an independent local model check before pulse optimization. This check used additional data within the same study, rather than independent-team reproduction. The scientific distinction is between fitting one spectrum and establishing a stable model suitable for subsequent control. More pulse optimization using unstable parameters would address a different uncertainty. Reuse of the adapted analysis on Q3 therefore concerned a procedure, with fresh spectra, references, and control synthesis still required.

The AI Theorist case changes explanatory assumptions rather than only a residual function or parameter-extraction procedure. For $\alpha$-RuCl$_3$, a constrained electronic starting state produced a strong optical feature absent from the observations. Releasing the orbital-occupation constraint reduced that feature and restored the observed dominant peak in the calculation. Updating quasiparticle energies while holding screening fixed changed the gap little; updating screening as well materially altered the predicted gap. Electron--hole interactions, treated through the Bethe--Salpeter equation, then produced candidate bound excitations below the calculated quasiparticle gap~\cite{asd1081}. These tests address different physical sources of disagreement: the electronic state, its screening response, and the interacting excited states. A shift of an optical peak is consequently not equivalent to a correction of the quasiparticle gap. The delegated decisions selected calculations that changed these physical assumptions, while human-acquired spectra and theoretical assessment constrained interpretation.

Agreement must also concern the observable actually measured. Reflection contrast, calculated dielectric response, and photocurrent are related but distinct quantities. Photocurrent requires charge generation and collection following excitation; matching excitation energies alone does not establish a quantitative account of that process. The study's proposed polarization contrast between candidate excitonic features supplies a further discriminating observation, but remained unmeasured~\cite{asd1081}. The model revisions therefore support a physically developed interpretation and testable predictions, while the uniqueness of the assignments and the complete photocurrent mechanism remain unresolved. Improving agreement with the observations used to guide revision cannot substitute for testing such additional consequences.

The contrast locates the delegated decision precisely: BEAR learns within researcher-designed response representations, the NMR agent revises parameter extraction with operator involvement, and the theoretical agents change explanatory calculations. A residual magnitude alone cannot select the appropriate response. Its interpretation depends on the diagnostics available to separate predictive error, measurement instability, and an inadequate physical explanation. The cases support these distinct scientific operations without demonstrating a system that autonomously controls all three.

